\subsection{The obstruction for every act of an eight-element monoid}

\begin{lemma}[A stable retract propagates Tor]
\label{C:stable-retract-Tor}
Let $R$ be a finite-dimensional algebra over a field, let $Q$ be
a finite-dimensional left module, and choose projective syzygies
$U=\Omega^2Q$ and $V=\Omega^{2+s}Q$, where $s\geq1$.
Suppose there are homomorphisms
\[
 F_i:U\longrightarrow V,\quad G_i:V\longrightarrow U,
 \qquad
 A_j:U\longrightarrow P_j,\quad B_j:P_j\longrightarrow U,
\]
with every $P_j$ projective, such that
\[
 1_U=\sum_i G_iF_i+\sum_j B_jA_j.                         \tag{\ref{C:stable-retract-Tor}.1}
\]
Then, for every right module $L$ and every $n\geq3$,
\[
 \operatorname{Tor}_n^R(L,Q)\ne0
 \quad\Longrightarrow\quad
 \operatorname{Tor}_{n+s}^R(L,Q)\ne0.
\]
If the first sum is empty, $U$ is projective and all these
degree-$n$ groups are zero.
\end{lemma}
\begin{proof}
Apply the additive functor $\operatorname{Tor}_q^R(L,-)$ for
$q\geq1$. Every composite through a $P_j$ becomes zero.
Consequently $\operatorname{Tor}_q^R(L,V)=0$ forces the identity
of $\operatorname{Tor}_q^R(L,U)$ to be zero. Dimension shifting,
with $q=n-2$, gives the implication. If the first sum is empty,
the identity of $U$ factors through $\bigoplus_jP_j$, so $U$
is a projective direct summand of that module.
\end{proof}

\begin{lemma}[A derived local Frobenius comparison]
\label{C:derived-Frobenius-Tor}
Let $R$ be a finite-dimensional $k$-algebra and let $A$ be a
quotient algebra which is local and Frobenius, with unique simple
module $k$ and $\operatorname{pd}_R A\leq1$ on the left.
Suppose a finite-dimensional left $R$-module $Q$ fits into
\[
 0\longrightarrow k\longrightarrow P\longrightarrow Q
 \longrightarrow0,
\]
where $P$ is projective and the first module is obtained by
restriction from the simple $A$-module. For every finite-dimensional
right $R$-module $L$, eventual vanishing of
$\operatorname{Tor}_n^R(L,Q)$ implies its vanishing for every
$n\geq3$.
\end{lemma}
\begin{proof}
We first record the needed fact about complexes over $A$.
Let $C$ be a derived $A$-module with finite-dimensional
cohomology concentrated in degrees zero and one. If
\[
 \operatorname{Ext}_A^n(k,C)=0\qquad\text{for all sufficiently large }n,
                                                               \tag{\ref{C:derived-Frobenius-Tor}.1}\label{eq:C-Frobenius-eventual-Ext}
\]
then these groups vanish for every $n\geq2$.

To prove this fact, choose a bounded-below injective resolution
$I^\bullet$ of $C$, starting in degree zero, with finite-dimensional
terms. Such a resolution is obtained by resolving a bounded
finite-dimensional representative, or by successive injective
envelopes. Since $A$ is local and Frobenius, its finite-dimensional
injectives are finite sums of the unique indecomposable injective
$A$. We may make the complex minimal directly. A unit matrix
entry in a differential can be isolated by invertible row and
column operations, splitting off a contractible complex
$A\xrightarrow{1}A$. The neighboring differential components
are zero by the complex equations. Remove these summands
successively by degree. Only finitely many can be removed at
each degree, so the resulting bounded-below complex is defined
degree by degree, remains an injective resolution of $C$, and
has every differential entry in $J=\operatorname{rad}(A)$.

These differentials kill the socles of the injective terms:
the socle of a local Frobenius algebra is annihilated on both
sides by its radical.
Thus the complex $\operatorname{Hom}_A(k,I^\bullet)$ has zero
differential. Each nonzero injective term has a nonzero socle,
so \eqref{eq:C-Frobenius-eventual-Ext} forces $I^n=0$ for all
sufficiently large $n$.
If the last nonzero term were $I^N$ with $N>1$, the equality
$H^N(C)=0$ would make $I^{N-1}\to I^N$ surjective. Since
$I^N$ is also projective, that surjection would split off
another contractible summand, contrary to minimality.
Equivalently, a matrix with entries in $J$ cannot surject onto
a nonzero free module. Therefore $I^n=0$ for every $n\geq2$,
proving the fact.

Now put $N=\operatorname{Hom}_k(L,k)$ and form the full derived
right adjoint
\[
 C=\operatorname{RHom}_R(A,N).
\]
The hypothesis on the left projective dimension of $A$ implies
that $H^j(C)=0$ outside $j=0,1$; its cohomology is finite-dimensional.
Restriction of modules along $R\to A$ is exact and has right
adjoint $\operatorname{Hom}_R(A,-)$. Applying this adjunction
to an injective resolution of $N$ gives
\[
 \operatorname{RHom}_R(k,N)
 \simeq
 \operatorname{RHom}_A\bigl(k,\operatorname{RHom}_R(A,N)\bigr).
                                                               \tag{\ref{C:derived-Frobenius-Tor}.2}\label{eq:C-Frobenius-derived-adjunction}
\]
Indeed, the right adjoint sends injective $R$-modules to
injective $A$-modules, so this is an equality of the appropriate
derived Hom computations. No ordinary change-of-rings Tor
formula is being assumed.

Duality of finite-dimensional vector spaces and the displayed
short exact sequence give
\[
 \operatorname{Hom}_k\bigl(\operatorname{Tor}_n^R(L,Q),k\bigr)
 \cong\operatorname{Ext}_R^n(Q,N)
 \cong\operatorname{Ext}_R^{n-1}(k,N)\qquad(n\geq2).
\]
Eventual vanishing therefore gives
\eqref{eq:C-Frobenius-eventual-Ext} for $C$ by
\eqref{eq:C-Frobenius-derived-adjunction}.
The proved minimal-injective-complex fact makes
$\operatorname{Ext}_R^m(k,N)=0$ for every $m\geq2$.
The same dimension shift and faithful vector-space duality
give $\operatorname{Tor}_n^R(L,Q)=0$ for all $n\geq3$.
\end{proof}

\begin{proposition}[No projective-plane act cone over eight elements]
\label{C:eight-act-cones}
For every finite monoid $M$ of order at most eight and every
nonempty finite right $M$-set $S$, the category $\mathcal C(M,S)$
of Lemma~\ref{C:act-cone-Tor} does not have classifying space
homotopy equivalent to $\mathbb CP^2$.
There is no bound on the cardinality of $S$.
The same exclusion holds for a triangular biset category with
other endomorphism monoid $N$ whenever the trivial left
$\mathbb QN^{\mathrm{op}}$-module is projective.
\end{proposition}
\begin{proof}
Orders at most seven are covered by
Proposition~\ref{C:seven-act-cones}. If
$H^2(BM;\mathbb Q)=0$, the relative-suspension argument in
that proposition makes every product of degree-two classes
zero. It remains to consider eight-element monoids with
nonzero rational second cohomology.

We use the complete own translation-table enumeration
\path{code/c_eight_H2_targets.cpp}, whose construction is
spelled out in the proof of Lemma~\ref{C:twelve-triangular}.
The minimal ideal must be a rectangular band of shape
$2\times2$, $2\times3$, or $3\times2$: fewer rows or columns
give group type, and a nontrivial local group with two rows
and columns would already require at least nine elements.
For the $2\times2$ shape, nonzero rational $H_2$ forces the
row and column augmentation characters of every extra element
to agree. The complete list of compatible translations and
all associative products gives 492 ordinary isomorphism
classes. The other two shapes give twelve classes each.
This produces the 516 tables indexed $0,\ldots,515$ in
\path{code/outputs/c_eight_H2_targets.tsv}; both orientations
are present. Exact rational balanced-module ranks show that
the six indices
\[
 496,498,499,506,512,513
\]
have $H_2(BM;\mathbb Q)=0$. Retaining them as harmless extra
entries, rather than dropping them, gives the following exact
partition of the entire 516-table list:
\[
\begin{array}{l|r}
 \text{certificate}&\text{number of tables}\\ \hline
 \text{original resolution of length at most two}&261\\
 \text{original two-periodic resolution tail}&157\\
 \text{split cancellation in an original length-three resolution}&5\\
 \text{additional projective second syzygy}&9\\
 \text{stable retract with recurrence step one}&65\\
 \text{stable retract with recurrence step two}&15\\
 \text{derived local Frobenius comparison}&4
\end{array}                                                     \tag{\ref{C:eight-act-cones}.1}\label{eq:C-eight-partition}
\]
The complete index-by-index partition, with the hash of every
underlying resolution certificate, is
\path{code/outputs/a_eight_cone_complete_partition.json}.
All resolutions and syzygies in this proof are for the actual
algebra $R=\mathbb QM$ and its trivial module $\mathbb Q$.
No equivalence of monoid classifying spaces is used to transfer
an arbitrary act.

For the first 261 cases, the program
\path{code/b_verify_resolutions.py} independently verifies the
complete integral projective resolution: each projective is
$\mathbb ZM e$ or $\mathbb ZM(e-f)$ for commuting idempotents
with $ef=fe=f$, each differential is equivariant, every image
is saturated, and the successive rank sums prove exactness.
The final differential is injective. In the 157 periodic cases
the same checker verifies the initial segment, equality of the
projectives from degree two onward, and the two zero matrix
products and rank sums that close the two-periodic tail.
Thus, for every right module $L$, either
$\operatorname{Tor}_n^R(L,\mathbb Q)=0$ for $n\geq3$, or
\[
 \operatorname{Tor}_{n+2}^R(L,\mathbb Q)
 \cong\operatorname{Tor}_n^R(L,\mathbb Q)\qquad(n\geq3).
\]
The exact indices, tables, verified ranks, and certificate
hashes are recorded in
\path{code/outputs/c_eight_cone_resolutions.json}.

The five split cancellations occur at indices
$182,307,382,407,437$. The last projective is $R\varepsilon$,
where $\varepsilon=g-f$, and its differential is literally
inclusion into a summand $R\eta$ of $P_2$, where $\eta=1-h$.
In the same index order, the triples $(g,f,h)$ are
\[
 (6,1,7),\ (7,2,6),\ (5,1,6),\ (6,3,7),\ (7,4,6).
\]
Direct multiplication gives
$\varepsilon^2=\varepsilon$, $\eta^2=\eta$, and
$\varepsilon\eta=\eta\varepsilon=\varepsilon$.
Hence $R\eta=R\varepsilon\oplus R(\eta-\varepsilon)$.
Cancel the contractible last summand to obtain a resolution
of length at most two. These integral ring equations, the
literal generator image, and the original complete resolutions
are independently checked by
\path{code/c_eight_split_cancellations.py}, with all five
records in its JSON output.

For the remaining syzygy calculations we use actual kernels
in the saved projective initial segments. The search
\path{code/a_eight_cone_syzygies.py} computes their full
$M$-actions and expresses $1_{\Omega^2\mathbb Q}$ as in
Lemma~\ref{C:stable-retract-Tor}, with all projective factors
equal to the free regular module $R$.
For the nine indices
\[
 32,137,371,450,459,463,464,497,501,
\]
the identity factors entirely through free modules, proving
that $\Omega^2\mathbb Q$ is projective.
For fifteen indices the later syzygy is $\Omega^4\mathbb Q$:
\[
 14,30,52,125,135,143,165,249,261,295,345,358,369,401,433.
\]
For the other 65 successful stable identities it is
$\Omega^3\mathbb Q$. Each certificate stores every intertwining
matrix, its rational coefficient, the two complete module
actions, and the exact identity obtained by summing the
composites. The independent verifier
\path{code/a_eight_cone_verify.py} imports neither the resolution
generator nor the syzygy search. It reconstructs each
projective as the image of its actual idempotent in $\mathbb QM$,
checks the kernel bases and exactness, and verifies every
intertwining equation and all 89 identities over $\mathbb Q$.
Consequently nonzero $\operatorname{Tor}_3^R(L,\mathbb Q)$
persists in arbitrarily high degrees in all 80 recurrence cases.

The four remaining indices are $9,10,16,18$.
Write $m_j$ for the element indexed $j$, so $m_0$ is the
global identity. Here $K=\{m_1,m_2,m_3,m_4\}$ is the minimal
$2\times2$ band, and
the quotient $A=\mathbb QM/\mathbb QK$ has basis $1,x,y,c$
with $c=m_5$, $x=m_6$, and $y=m_7$ in the table coordinates.
Its multiplication is exactly that of
Lemma~\ref{B:binary-pencil-Tor}, with respective matrices
\[
 \begin{pmatrix}0&1\\1&0\end{pmatrix},\quad
 \begin{pmatrix}0&1\\1&1\end{pmatrix},\quad
 \begin{pmatrix}1&0\\0&1\end{pmatrix},\quad
 \begin{pmatrix}1&0\\1&1\end{pmatrix}.
\]
Each is invertible, so the coefficient-of-$c$ pairing proves
that $A$ is local Frobenius. The ideal $\mathbb QK$ is the
direct sum of its two column modules, $Rm_1$ and $Rm_2$.
Both are projective, so
$\operatorname{pd}_R A\leq1$.
The augmentation $Rm_1\to\mathbb Q$ has kernel spanned by
$m_3-m_1$. Every nonidentity element kills this vector, whereas
the global identity, indexed zero, fixes it. Thus
$\Omega_R\mathbb Q$ is precisely the simple $A$-module.
All these assertions are direct multiplication-table equations
and are separately checked for each of the four records in
\path{code/outputs/a_eight_cone_verified.json}.
Lemma~\ref{C:derived-Frobenius-Tor} therefore says that eventual
Tor vanishing forces $\operatorname{Tor}_3^R(L,\mathbb Q)=0$
in these last cases as well.

The rows of \eqref{eq:C-eight-partition} now show, for every
finite-dimensional right
$R$-module $L$, that nonzero degree-three Tor is incompatible
with bounded Tor. Apply Lemma~\ref{C:act-cone-Tor} with
$L=\ker(\mathbb QS\to\mathbb Q)$. A projective-plane act cone
would have nonzero rational fourth homology and zero homology
in every higher degree, which is impossible.
The biset statement follows by the exact second-vertex
coinvariant argument in that lemma.
\end{proof}
